% 附录 A：附件与数据说明
% ---------------------------------------------------------------------------
% 写作阶段填写。列出附件清单、字段含义、清洗口径，以及每个代码文件的作用，
% 让评阅人能根据附录定位正文的每个结果来自哪个附件与脚本。
% ---------------------------------------------------------------------------
%
% 建议结构（示例，按本题实际情况调整）：
%
% \begin{description}
%   \item[附件一] 男胎检测数据（1082 行 / 267 名孕妇）：孕周、BMI、Y 染色体浓度等。
%   \item[附件二] 女胎检测数据（605 行 / 147 名孕妇）：X 染色体浓度、Z 值、AB 列标注等。
% \end{description}
%
% 数据清洗口径：孕周"11w+6"解析为 11.857 周；同一孕妇多次采血按孕妇代码分组；
% 女胎以 AB 列非整倍体为判定金标准，AE 列（出生结局）不用于监督训练。
%
% 代码文件说明：
% \begin{table}[htbp]
%   \centering
%   \begin{tabular}{ll}
%     \toprule
%     文件 & 作用 \\
%     \midrule
%     code/load\_data.py & 数据读取与孕周解析 \\
%     code/q1\_model.py  & 问题一关系模型 \\
%     % ... 按实际代码文件填写
%     \bottomrule
%   \end{tabular}
% \end{table}
